\providecommand{\flag}[1]{{\color{red}#1}}  % weakness statements: author to soften or remove
\section{Validating Our Model}
\label{sec:results}
\label{sec:res-table}

\paraq{What does this section check?}
The floorline of Section~\ref{sec:roofline} is correct when the strategy it picks is the strategy
that measures fastest. We check this on the grid of Section~\ref{sec:hybrids}, reusing its
measurements: for each configuration, the floorline picks the strategy with the smallest modelled
overhead among DGL, gsplit, and the hybrids $0<k<L$ (Section~\ref{sec:cm-layers}), and we compare
that pick with the fastest measured strategy by forward plus backward time, the part of an epoch
the floorline models. No measured time is fitted.
% Detail: the pipeline is not a candidate in Table~\ref{tab:final-grid}; Section~\ref{sec:hybrids}
% measured it slower than baseline gsplit in all 40 configurations.

\paraq{How is the check scored?}
Every strategy is a split point $k$, the number of leading data-parallel layers: gsplit is $k{=}0$,
the hybrids are $0<k<L$, and DGL is $k{=}L$. A pick \emph{agrees} with the measured winner when
its split point is the winner's or one layer from it; a pick two or more layers away is a miss.
The floorline ranks strategies by a lower bound on their overhead; it does not predict their times,
so we score the distance in layers and not in seconds.

\paraq{Does the floorline pick the winner?}
Table~\ref{tab:final-grid} is the validation. It has the layout of
Tables~\ref{tab:sec1-products} and~\ref{tab:sec1-papers}: a hidden sweep at $L{=}3$ and a depth
sweep at $h{=}1024$ per model, for both graphs and both interconnects. Each cell is the floorline's
pick. Dark green is an exact match; light green is a pick one layer from the measured winner, which
is named in parentheses; red is a miss.
% ---------------------------------------------------------------------------------------------
% COUNTS (for review; not printed) -- from gen_roofline_grid.py's console summary, 2026-11-18.
%   Section 2 grid, 40 distinct settings (the default L=3,h=1024 is printed twice in the table):
%     exact 20, one layer off 12, miss (two or more layers) 8  ->  agree 32 / 40.
%     exact only:  PCIe 14/20, NVLink 6/20.
%     products SAGE  PCIe 4/5, NVLink 2/5      papers100M SAGE  PCIe 3/5, NVLink 3/5
%     products GAT   PCIe 3/5, NVLink 1/5      papers100M GAT   PCIe 4/5, NVLink 0/5
%     one layer off (12), pick -> measured winner:
%       DGL -> h2 at L=3 (6), DGL -> h3 at L=4 (1), h1 -> DGL at L=2 (2),
%       h1 -> gsplit (1), h1 -> h2 (2)
%     misses (8), all NVLink, all with k_pick < k_best:
%       gsplit -> h2 at L=3 (4), gsplit -> h3 at L=4 (2), h1 -> DGL at L=3 (1), h1 -> h3 at L=4 (1)
%   Extended sweep (69 settings, exact match 68%): results/roofline/agreement.csv,
%   tables/roofline_agreement{,_summary}.tex, tables/roofline_grid_all.tex -- not printed.
% ---------------------------------------------------------------------------------------------

\begin{table}[h]
\centering
\caption{Floorline pick vs.\ measured winner on the grid of Section~\ref{sec:hybrids}. Each cell is
the pick (D~=~DGL, G~=~gsplit, h$k$~=~hybrid with split point $k$). Dark green~=~exact match;
light green~=~one layer from the measured winner (in parentheses); red~=~two or more layers away. Two single-axis
ablations per model crossing at the default $(L{=}3,\,h{=}1024)$: a hidden sweep at $L{=}3$, then a
depth sweep at $h{=}1024$. Four A100 GPUs, global mini-batch $1024$, forward plus backward time.
GAT shuffle time is scaled by the calibration constant of Section~\ref{sec:cm-estimate}.}
\label{tab:final-grid}
% generated by gen_roofline_grid.py -- do not edit
\begin{tabular}{|l|c|c|c|c|c|c|}
\hline
\multicolumn{3}{|c|}{} & \multicolumn{2}{c|}{\textbf{products}} & \multicolumn{2}{c|}{\textbf{papers100M}} \\
\cline{4-7}
\textbf{Model} & $\mathbf{h}$ & $\mathbf{L}$ & PCIe (emul.) & NVLink & PCIe (emul.) & NVLink \\
\hline
\multirow{6}{*}{\textbf{SAGE}} & $256$ & \multirow{3}{*}{$3$} & \cellcolor{green!45}D & \cellcolor{green!45}D & \cellcolor{green!15}D\,(h2) & \cellcolor{red!30}G\,(h2) \\
 & $1024$ &  & \cellcolor{green!45}D & \cellcolor{green!15}D\,(h2) & \cellcolor{green!45}D & \cellcolor{green!45}h2 \\
 & $4096$ &  & \cellcolor{green!15}D\,(h2) & \cellcolor{green!15}D\,(h2) & \cellcolor{green!15}D\,(h2) & \cellcolor{green!45}h2 \\
\cline{2-7}
 & \multirow{3}{*}{$1024$} & $2$ & \cellcolor{green!45}D & \cellcolor{green!45}D & \cellcolor{green!45}D & \cellcolor{green!45}D \\
 &  & $3$ & \cellcolor{green!45}D & \cellcolor{green!15}D\,(h2) & \cellcolor{green!45}D & \cellcolor{green!45}h2 \\
 &  & $4$ & \cellcolor{green!45}D & \cellcolor{red!30}h1\,(h3) & \cellcolor{green!45}h3 & \cellcolor{red!30}G\,(h3) \\
\hline
\multirow{6}{*}{\textbf{GAT}} & $256$ & \multirow{3}{*}{$3$} & \cellcolor{green!45}D & \cellcolor{red!30}G\,(h2) & \cellcolor{green!45}D & \cellcolor{red!30}h1\,(D) \\
 & $1024$ &  & \cellcolor{green!15}D\,(h2) & \cellcolor{red!30}G\,(h2) & \cellcolor{green!45}h2 & \cellcolor{green!15}h1\,(h2) \\
 & $4096$ &  & \cellcolor{green!45}h2 & \cellcolor{red!30}G\,(h2) & \cellcolor{green!45}h2 & \cellcolor{green!15}h1\,(h2) \\
\cline{2-7}
 & \multirow{3}{*}{$1024$} & $2$ & \cellcolor{green!45}D & \cellcolor{green!45}D & \cellcolor{green!15}h1\,(D) & \cellcolor{green!15}h1\,(D) \\
 &  & $3$ & \cellcolor{green!15}D\,(h2) & \cellcolor{red!30}G\,(h2) & \cellcolor{green!45}h2 & \cellcolor{green!15}h1\,(h2) \\
 &  & $4$ & \cellcolor{green!15}D\,(h3) & \cellcolor{green!15}h1\,(G) & \cellcolor{green!45}h3 & \cellcolor{red!30}G\,(h3) \\
\hline
\end{tabular}

\end{table}

\paraq{Where does it agree, and where does it miss?}
\expected{to verify against the table.} Every PCIe cell is green. The red cells are all on NVLink,
and in each of them the floorline picks fewer data-parallel layers than the measured winner has:
it picks gsplit or $k{=}1$ where a hybrid with a larger $k$, or DGL, measures fastest. The
light-green cells are mostly DGL picked where the hybrid $k{=}L{-}1$ wins.
Most red cells are GAT: its shuffle moves narrow rows of the transformed features, whose
small-message all-to-all runs below peak bandwidth, and the single calibration constant of
Section~\ref{sec:cm-estimate} absorbs only the average of that inefficiency.

\paraq{What is outside this check?}
Two things. First, the pipeline: the floorline credits it with the perfect overlap it can reach in
principle, whereas our implementation synchronises the two micro-batches at every step and
serialises their shuffles on one communicator, so the measured pipeline never wins
(Section~\ref{sec:hybrids}) and it is not a candidate in Table~\ref{tab:final-grid}. Second,
sampling and feature loading: they are part of the epoch but not of the floorline, so the check
uses forward plus backward time only.

\paraq{What does this section conclude?}
The floorline picks the measured winner, or a split point one layer from it, in $32$ of the $40$
settings, and in every PCIe setting. The $8$ misses are all on NVLink, where it picks fewer
data-parallel layers than the measured winner has. Built from one sampling pass and three measured
rates, with no measured time fitted, the model selects the strategy without the sweep of
Section~\ref{sec:hybrids}.
